% =====================================================================
% Chủ đề 5 — Tư duy logic và Câu đố — phần a: lý thuyết, ví dụ, bài 1–13
% =====================================================================
% helper macros (prefix \lga..., guarded by \providecommand)
% \lgaScale{left}{right}: chiếc cân hai đĩa đang thăng bằng
\providecommand{\lgaScale}[2]{\begin{tikzpicture}[baseline=-0.3cm,font=\small]
  \draw[Ink,line width=1.1pt] (-2.25,0)--(-0.45,0) (0.45,0)--(2.25,0);
  \draw[Ink,line width=1pt] (-1.35,0)--(-1.35,-0.5)--(1.35,-0.5)--(1.35,0);
  \draw[Ink,line width=1pt] (-0.17,-0.78)--(0,-0.5)--(0.17,-0.78);
  \node[above=0.5pt,inner sep=1pt] at (-1.35,0){#1};\node[above=0.5pt,inner sep=1pt] at (1.35,0){#2};\end{tikzpicture}}
% hình dùng trên cân
\providecommand{\lgaSq}{\tikz[baseline=0pt]\draw[fill=tmBlue!65,draw=tmBlue!55!black,line width=0.6pt] (0,0) rectangle (0.3,0.3);}
\providecommand{\lgaTri}{\tikz[baseline=0pt]\draw[fill=tmOrange!85,draw=tmOrange!55!black,line width=0.6pt] (0,0)--(0.34,0)--(0.17,0.31)--cycle;}
\providecommand{\lgaCir}{\tikz[baseline=0pt]\draw[fill=tmGreen!65,draw=tmGreen!45!black,line width=0.6pt] (0.15,0.15) circle (0.15);}
% trái cây: táo (đỏ, có cuống), cam (cam), xoài (xanh, bầu dục)
\providecommand{\lgaApple}{\tikz[baseline=0pt]{\draw[fill=AnsRed!80,draw=AnsRed!55!black,line width=0.5pt] (0.15,0.15) circle (0.15);
  \draw[tmGreen!50!black,line width=0.8pt] (0.15,0.29)--(0.2,0.39);}}
\providecommand{\lgaOrg}{\tikz[baseline=0pt]{\draw[fill=tmOrange,draw=tmOrange!60!black,line width=0.5pt] (0.17,0.17) circle (0.17);
  \fill[tmOrange!45!black] (0.17,0.3) circle (0.025);}}
\providecommand{\lgaMango}{\tikz[baseline=0pt]{\draw[fill=tmGreen!70,draw=tmGreen!45!black,line width=0.5pt] (0.27,0.15) ellipse (0.27 and 0.15);}}
% hộp màu (Bài 23)
\providecommand{\lgaBox}[1]{\tikz[baseline=0pt]\draw[fill=#1!70,draw=#1!55!black,line width=0.6pt,rounded corners=1pt] (0,0) rectangle (0.34,0.3);}
% \lgaMagic{a}{b}{c}{d}{e}{f}{g}{h}{i}: lưới 3x3 (đọc theo hàng)
\providecommand{\lgaMagic}[9]{\begin{tikzpicture}[baseline=(current bounding box.center),font=\normalsize]
  \fill[tmBlue!5] (0,0) rectangle (2.7,2.7);
  \draw[Navy,line width=0.9pt] (0,0) grid[step=0.9] (2.7,2.7);
  \node at (0.45,2.25){#1};\node at (1.35,2.25){#2};\node at (2.25,2.25){#3};
  \node at (0.45,1.35){#4};\node at (1.35,1.35){#5};\node at (2.25,1.35){#6};
  \node at (0.45,0.45){#7};\node at (1.35,0.45){#8};\node at (2.25,0.45){#9};\end{tikzpicture}}
% \lgaFour{c/r/giá trị,...}: lưới 4x4 chia 4 ô 2x2 (c = cột 1..4 từ trái, r = hàng 1..4 từ trên)
\providecommand{\lgaFour}[1]{\begin{tikzpicture}[baseline=(current bounding box.center),font=\normalsize]
  \fill[tmBlue!4] (0,0) rectangle (2.8,2.8);
  \draw[Navy!60,line width=0.5pt] (0,0) grid[step=0.7] (2.8,2.8);
  \draw[Navy,line width=1.6pt] (0,0) rectangle (2.8,2.8) (1.4,0)--(1.4,2.8) (0,1.4)--(2.8,1.4);
  \foreach \c/\r/\v in {#1}{\node at ({0.7*\c-0.35},{2.8-0.7*\r+0.35}){\v};}\end{tikzpicture}}
% \lgaPips{n}: các chấm của mặt xúc xắc trong hình vuông đơn vị
\providecommand{\lgaPips}[1]{\ifcase#1\or
  \fill (0.5,0.5) circle (0.085);\or
  \fill (0.27,0.73) circle (0.085) (0.73,0.27) circle (0.085);\or
  \fill (0.25,0.75) circle (0.085) (0.5,0.5) circle (0.085) (0.75,0.25) circle (0.085);\or
  \fill (0.27,0.27) circle (0.085) (0.73,0.73) circle (0.085) (0.27,0.73) circle (0.085) (0.73,0.27) circle (0.085);\or
  \fill (0.25,0.25) circle (0.085) (0.75,0.75) circle (0.085) (0.25,0.75) circle (0.085) (0.75,0.25) circle (0.085) (0.5,0.5) circle (0.085);\or
  \fill (0.27,0.22) circle (0.08) (0.73,0.22) circle (0.08) (0.27,0.5) circle (0.08) (0.73,0.5) circle (0.08) (0.27,0.78) circle (0.08) (0.73,0.78) circle (0.08);\fi}
% \lgaCube{trên}{trước}{phải}: con xúc xắc nhìn nghiêng (thấy 3 mặt)
\providecommand{\lgaCube}[3]{\begin{tikzpicture}[baseline=0.4cm,scale=0.85,line join=round]
  \draw[Ink,line width=0.9pt,fill=white] (0,0) rectangle (1,1);
  \draw[Ink,line width=0.9pt,fill=white] (0,1)--(1,1)--(1.45,1.45)--(0.45,1.45)--cycle;
  \draw[Ink,line width=0.9pt,fill=white] (1,0)--(1.45,0.45)--(1.45,1.45)--(1,1)--cycle;
  \begin{scope}[Ink]\lgaPips{#2}\end{scope}
  \begin{scope}[Ink,cm={1,0,0.45,0.45,(0,1)}]\lgaPips{#1}\end{scope}
  \begin{scope}[Ink,cm={0.45,0.45,0,1,(1,0)}]\lgaPips{#3}\end{scope}\end{tikzpicture}}
% \lgaTriEight{đỉnh}{t1}{t2}{góc trái}{p1}{p2}{góc phải}{đáy giữa}: tam giác 8 vòng tròn
\providecommand{\lgaTriEight}[8]{\begin{tikzpicture}[baseline=(current bounding box.center),font=\small]
  \draw[Navy,line width=0.9pt] (0,2.7)--(-1.5,0)--(1.5,0)--cycle;
  \foreach \x/\y in {0/2.7,-0.5/1.8,-1/0.9,-1.5/0,0.5/1.8,1/0.9,1.5/0,0/0}{\draw[Navy,fill=white,line width=0.9pt] (\x,\y) circle (0.27);}
  \node at (0,2.7){#1};\node at (-0.5,1.8){#2};\node at (-1,0.9){#3};\node at (-1.5,0){#4};
  \node at (0.5,1.8){#5};\node at (1,0.9){#6};\node at (1.5,0){#7};\node at (0,0){#8};\end{tikzpicture}}
% \lgaTriSix{đỉnh}{góc trái}{góc phải}{giữa cạnh trái}{giữa đáy}{giữa cạnh phải}: tam giác 6 vòng tròn
\providecommand{\lgaTriSix}[6]{\begin{tikzpicture}[baseline=(current bounding box.center),font=\small]
  \draw[Navy,line width=0.9pt] (0,2.2)--(-1.27,0)--(1.27,0)--cycle;
  \foreach \x/\y in {0/2.2,-1.27/0,1.27/0,-0.635/1.1,0/0,0.635/1.1}{\draw[Navy,fill=white,line width=0.9pt] (\x,\y) circle (0.27);}
  \node at (0,2.2){#1};\node at (-1.27,0){#2};\node at (1.27,0){#3};
  \node at (-0.635,1.1){#4};\node at (0,0){#5};\node at (0.635,1.1){#6};\end{tikzpicture}}
% \lgaSqEight{TL}{trên}{TR}{phải}{BR}{dưới}{BL}{trái}: hình vuông 8 vòng tròn (theo chiều kim đồng hồ)
\providecommand{\lgaSqEight}[8]{\begin{tikzpicture}[baseline=(current bounding box.center),font=\small]
  \draw[Navy,line width=0.9pt] (0,0) rectangle (2.4,2.4);
  \foreach \x/\y in {0/2.4,1.2/2.4,2.4/2.4,2.4/1.2,2.4/0,1.2/0,0/0,0/1.2}{\draw[Navy,fill=white,line width=0.9pt] (\x,\y) circle (0.27);}
  \node at (0,2.4){#1};\node at (1.2,2.4){#2};\node at (2.4,2.4){#3};\node at (2.4,1.2){#4};
  \node at (2.4,0){#5};\node at (1.2,0){#6};\node at (0,0){#7};\node at (0,1.2){#8};\end{tikzpicture}}
% \lgaCross{trên}{trái}{giữa}{phải}{dưới}: hình chữ thập 5 ô
\providecommand{\lgaCross}[5]{\begin{tikzpicture}[baseline=(current bounding box.center),font=\normalsize]
  \foreach \x/\y in {1/2,0/1,1/1,2/1,1/0}{\draw[Navy,fill=tmBlue!5,line width=0.9pt] (\x*0.8,\y*0.8) rectangle ++(0.8,0.8);}
  \node at (1.2,2.0){#1};\node at (0.4,1.2){#2};\node at (1.2,1.2){#3};\node at (2.0,1.2){#4};\node at (1.2,0.4){#5};\end{tikzpicture}}
% \lgaWheel{k}: 26 chữ cái xếp vòng theo chiều kim đồng hồ, mũi tên minh hoạ phép dịch k bước từ A
\providecommand{\lgaWheel}[1]{\begin{tikzpicture}[baseline=(current bounding box.center),font=\scriptsize\sffamily]
  \draw[Navy!50,line width=0.8pt] (0,0) circle (1.95);
  \foreach \l [count=\i from 0] in {A,B,C,D,E,F,G,H,I,J,K,L,M,N,O,P,Q,R,S,T,U,V,W,X,Y,Z}{
    \node[circle,fill=white,draw=Navy!60,inner sep=0.6pt,minimum size=0.36cm] at ({90-\i*360/26}:1.95){\l};}
  \draw[->,tmOrange,line width=1pt] ({90-3}:1.4) arc ({90-3}:{90-#1*360/26+3}:1.4);
  \node[font=\footnotesize\bfseries,text=tmOrange!80!black] at (0,0){$+#1$};\end{tikzpicture}}

\begin{knowbox}
\textbf{Ý tưởng lớn.} Tư duy logic là \emph{nghệ thuật rút ra điều chắc chắn từ những điều đã biết}. Một câu đố logic
thường không cần phép tính phức tạp: cái khó là phải \emph{đọc kĩ}, \emph{ghi chép có hệ thống} và \emph{không bỏ sót trường hợp nào}.
Trong chủ đề này, em sẽ làm quen với các ``dụng cụ'' sau:
\begin{itemize}[leftmargin=1.4em,itemsep=1pt,topsep=2pt]
  \item \textbf{câu nói thật -- nói dối}, câu ``đúng một nửa'', và cách phủ định một câu;
  \item \textbf{xét trường hợp} và \textbf{phản chứng} (giả sử rồi tìm mâu thuẫn);
  \item \textbf{sắp thứ tự} bằng sơ đồ mũi tên, trục số; phân biệt ``\emph{chắc chắn đúng}'' với ``\emph{có thể đúng}'';
  \item \textbf{mật mã dịch vòng} (mật mã Caesar) và \textbf{số dư} khi chia cho $26$ (hay cho $7$ với các ngày trong tuần);
  \item \textbf{ma phương}, \textbf{chiếc cân} (chính là phương trình), \textbf{tổng trên hình} và \textbf{xúc xắc};
  \item \textbf{hình vuông Latin} và \textbf{sudoku} $4\times 4$.
\end{itemize}
Bốn thói quen vàng: \emph{(1) viết ra -- kí hiệu hoá; (2) vẽ ra -- sơ đồ, bảng; (3) xét đủ mọi trường hợp; (4) thử lại
đáp án với \textbf{tất cả} điều kiện của đề.}
\end{knowbox}
\vspace{0.3cm}

% ---------------------------------------------------------------------
\needspace{9\baselineskip}\shead{Câu đúng, câu sai: người nói thật và người nói dối}
\begin{knowbox}
\textbf{1. Mệnh đề.} Một câu khẳng định hoặc \emph{đúng}, hoặc \emph{sai} (không thể vừa đúng vừa sai) gọi là một mệnh đề.
Ví dụ: ``$7$ là số lẻ'' (đúng), ``Hà Nội ở phía nam TP.~Hồ Chí Minh'' (sai).

\textbf{2. Người nói thật và người nói dối.} Trong các câu đố, \emph{người nói thật} thì mọi câu họ nói đều đúng;
\emph{người nói dối} thì mọi câu họ nói đều sai. Vì vậy:
\begin{itemize}[leftmargin=1.4em,itemsep=1pt,topsep=2pt]
  \item nếu biết người nói là ai, ta biết câu nói đúng hay sai;
  \item ngược lại, nếu biết câu nói đúng hay sai, ta biết người nói là ai.
\end{itemize}

\textbf{3. Phủ định một câu (nói ``ngược lại'').}
\begin{center}
\renewcommand{\arraystretch}{1.25}
\begin{tabular}{|l|l|}
\hline
\rowcolor{tmBlue!10} Câu & Câu phủ định (đúng khi câu kia sai) \\ \hline
An \emph{và} Bình đều nói dối & \emph{Ít nhất một} trong hai bạn nói thật \\ \hline
\emph{Mọi} quả táo đều đỏ & \emph{Có ít nhất một} quả táo không đỏ \\ \hline
Hoa là hoa hồng đỏ & Hoa không đỏ, \emph{hoặc} không phải hoa hồng (hoặc cả hai) \\ \hline
\end{tabular}
\end{center}

\textbf{4. Câu ``đúng một nửa''.} Nếu một câu gồm hai ý và ta biết \emph{đúng một} ý đúng, thì:
ý thứ nhất đúng $\Rightarrow$ ý thứ hai sai; \ ý thứ nhất sai $\Rightarrow$ ý thứ hai đúng. Chỉ cần đoán đúng/sai cho \emph{một} ý,
mọi ý còn lại sẽ ``kéo theo'' như một dây domino.

\textbf{5. Nói dối theo ngày.} Một người ``nói dối vào thứ Ba'' thì mọi câu họ nói vào thứ Ba đều sai. Muốn biết ngày nào họ
\emph{có thể} nói một câu, ta lập bảng: mỗi ngày, xem câu đó đúng hay sai, rồi so với việc hôm đó họ nói thật hay dối.
\end{knowbox}
\begin{proofbox}
Vì sao không ai có thể nói ``Tôi là người nói dối''? Nếu người nói là người nói thật thì câu đó phải đúng, tức là họ nói dối --
vô lí. Nếu người nói là người nói dối thì câu đó phải sai, tức là họ \emph{không} nói dối -- cũng vô lí. Cả hai trường hợp đều
mâu thuẫn, nên câu này không thể được nói ra (trên ``hòn đảo'' chỉ có người nói thật và người nói dối).
\end{proofbox}
\begin{warnbox}
Phủ định của ``An và Bình đều nói dối'' \emph{không phải} là ``An và Bình đều nói thật''. Chỉ cần \emph{một} bạn nói thật là
câu ban đầu đã sai rồi.
\end{warnbox}
\begin{tipbox}
Gặp câu đố thật -- dối, hãy kẻ một bảng nhỏ: mỗi hàng là một khả năng (ai thật, ai dối), mỗi cột là một câu nói. Điền Đ/S
vào từng ô; hàng nào không mâu thuẫn chính là đáp án.
\end{tipbox}

% ---------------------------------------------------------------------
\needspace{9\baselineskip}\shead{Xét trường hợp và phản chứng}
\begin{knowbox}
\textbf{1. Xét trường hợp.} Khi chưa biết điều gì đó, ta chia ra \emph{tất cả} các khả năng có thể xảy ra (không bỏ sót),
rồi xét lần lượt từng khả năng. Khả năng nào dẫn đến điều vô lí (mâu thuẫn với đề bài) thì bị loại. Khả năng còn lại
chính là đáp án.

\textbf{2. Phản chứng.} Muốn chứng tỏ một điều $P$ đúng, ta \emph{giả sử} $P$ sai, rồi suy luận tiếp. Nếu đi đến một điều
vô lí thì giả sử là sai, vậy $P$ đúng.

\textbf{3. Hai câu mâu thuẫn.} Nếu hai câu không thể cùng đúng và cũng không thể cùng sai (chẳng hạn Long nói ``Hà nói dối''
về câu nói của Hà) thì \emph{đúng một} trong hai câu đúng. Đây là ``chìa khoá'' của rất nhiều câu đố: biết ngay được một câu đúng,
một câu sai trong cặp đó.

\textbf{4. Xếp chỗ quanh bàn tròn.} Khi xoay cả bàn, ai ngồi cạnh ai không thay đổi. Vì vậy ta có thể \emph{cố định} chỗ ngồi
của một người, rồi chỉ xét chỗ ngồi của những người còn lại. Quanh bàn tròn $5$ chỗ, mỗi người có đúng $2$ người ngồi cạnh,
và đúng $2$ người \emph{không} ngồi cạnh.
\end{knowbox}
\begin{figC}
  \foreach \i in {0,...,4}{\coordinate (P\i) at ({90-\i*72}:1.35);}
  \draw[Navy!40,line width=5pt] (0,0) circle (0.75);
  \foreach \i/\n in {0/1,1/2,2/3,3/4,4/5}{\node[circle,draw=Navy,fill=tmBlue!8,line width=0.8pt,minimum size=0.55cm] at (P\i){\n};}
  \draw[tmGreen!70!black,line width=1pt] (P0) to[bend left=25] (P1);
  \draw[tmGreen!70!black,line width=1pt] (P0) to[bend right=25] (P4);
  \draw[AnsRed!70,dashed,line width=1pt] (P0)--(P2) (P0)--(P3);
  \node[align=left,font=\footnotesize,text=Slate,anchor=west] at (1.6,0.3){\textcolor{tmGreen!60!black}{liền nét}: ghế $1$ ngồi cạnh ghế $2$ và ghế $5$};
  \node[align=left,font=\footnotesize,text=Slate,anchor=west] at (1.6,-0.3){\textcolor{AnsRed!80}{nét đứt}: ghế $1$ không ngồi cạnh ghế $3$, ghế $4$};
\end{figC}
\begin{proofbox}
Vì sao phản chứng là một lập luận đúng? Một mệnh đề chỉ có thể đúng hoặc sai. Nếu ``$P$ sai'' dẫn đến điều vô lí thì khả năng
``$P$ sai'' không thể xảy ra. Chỉ còn lại khả năng ``$P$ đúng''. Đây chính là phép xét trường hợp với hai trường hợp: $P$ đúng
hoặc $P$ sai.
\end{proofbox}
\begin{tipbox}
Hãy bắt đầu từ thông tin ``chặt'' nhất -- thông tin cho ít khả năng nhất. Chẳng hạn ``Giang và Khoa \emph{không} ngồi cạnh nhau''
quanh bàn $5$ chỗ nghĩa là giữa họ có \emph{đúng một} người ở một phía. Ít khả năng thì ít phải xét!
\end{tipbox}

% ---------------------------------------------------------------------
\needspace{9\baselineskip}\shead{Sắp thứ tự; ``chắc chắn đúng'' và ``có thể đúng''}
\begin{knowbox}
\textbf{1. Kí hiệu hoá.} Dịch mỗi câu thành một bất đẳng thức: ``Minh chạy nhanh hơn Hùng'' viết $M > H$;
``Kiều ăn ít hơn Long'' viết $K < L$ (hay $L > K$).

\textbf{2. Tính bắc cầu.} Nếu $A > B$ và $B > C$ thì $A > C$. Nhờ đó ta nối các bất đẳng thức thành một dãy.

\textbf{3. Sơ đồ.} Vẽ mỗi người là một điểm trên một \emph{trục dọc}: ai ``hơn'' thì đặt cao hơn; hoặc vẽ mũi tên
từ người ``kém'' đến người ``hơn''. Hai người không nối được với nhau bằng một dãy mũi tên cùng chiều thì \emph{chưa so sánh được}.

\textbf{4. Thứ tự trong hàng.} Nếu một bạn đứng thứ $k$ tính từ đầu hàng và thứ $m$ tính từ cuối hàng thì hàng có
$k + m - 1$ người.

\textbf{5. ``Chắc chắn đúng'' và ``có thể đúng''.}
\begin{itemize}[leftmargin=1.4em,itemsep=1pt,topsep=2pt]
  \item Câu \textbf{chắc chắn đúng} (luôn đúng): đúng trong \emph{mọi} trường hợp thoả mãn đề bài.
  \item Câu \textbf{có thể đúng}: đúng trong \emph{ít nhất một} trường hợp thoả mãn đề bài.
  \item Muốn bác bỏ một câu ``chắc chắn đúng'', chỉ cần \emph{một phản ví dụ}: một trường hợp thoả đề mà câu đó sai.
\end{itemize}

\textbf{6. ``Mọi A đều là B''.} Từ câu này ta suy ra được: ``ai \emph{không} là B thì \emph{không} là A''. Nhưng ta
\emph{không} suy ra được ``ai là B thì là A''.
\end{knowbox}
\begin{figC}
  \begin{scope}
  \foreach \y/\n in {2.4/Quân,1.8/Tâm,1.2/Minh,0.6/Phúc,0/Hùng}{\fill[tmBlue] (0,\y) circle (2.2pt); \node[right,font=\small] at (0.12,\y){\n};}
  \draw[->,Navy,line width=0.8pt] (-0.35,-0.15)--(-0.35,2.6) node[above,font=\scriptsize,text=Slate]{nhanh hơn};
  \end{scope}
  \begin{scope}[xshift=5.2cm]
  \draw[tmBlue,line width=1pt,fill=tmBlue!8] (0.9,1.1) ellipse (1.7 and 1.15);
  \draw[tmOrange,line width=1pt,fill=tmOrange!15] (0.55,1.0) ellipse (0.75 and 0.6);
  \node[font=\scriptsize,text=tmOrange!70!black] at (0.55,1.0){CLB cờ vua};
  \node[font=\scriptsize,text=tmBlue!80!black] at (1.75,1.85){biết bơi};
  \fill[AnsRed] (3.2,1.1) circle (2pt) node[right,font=\scriptsize]{không biết bơi};
  \end{scope}
\end{figC}
\begin{proofbox}
Vì sao đứng thứ $k$ từ đầu và thứ $m$ từ cuối thì hàng có $k+m-1$ người? Đếm từ đầu đến bạn ấy được $k$ người, đếm từ bạn ấy
đến cuối được $m$ người. Cộng lại $k+m$, nhưng bạn ấy được đếm \emph{hai lần}, nên phải trừ $1$.

Vì sao ``mọi A là B'' cho ta ``không là B thì không là A''? Nhìn hình bên phải: vòng ``CLB cờ vua'' nằm \emph{trọn} trong vòng
``biết bơi''. Một điểm nằm ngoài vòng lớn thì chắc chắn nằm ngoài vòng nhỏ. Còn một điểm trong vòng lớn có thể nằm trong
hoặc ngoài vòng nhỏ.
\end{proofbox}
\begin{warnbox}
Thấy ``$A > B$'' và ``$A > C$'' đừng vội kết luận gì về $B$ và $C$: chúng \emph{chưa so sánh được}. Một câu như ``$B > C$''
khi đó chỉ \emph{có thể} đúng, không \emph{chắc chắn} đúng.
\end{warnbox}

% ---------------------------------------------------------------------
\needspace{9\baselineskip}\shead{Mật mã dịch vòng và số dư}
\begin{knowbox}
\textbf{1. Mật mã Caesar.} Xếp $26$ chữ cái tiếng Anh A, B, C, \ldots, Z theo vòng tròn (sau Z lại đến A). Mã hoá với
\emph{bước dịch} $k$: thay mỗi chữ bằng chữ đứng sau nó $k$ vị trí theo chiều kim đồng hồ. Giải mã: dịch \emph{ngược} lại $k$ vị trí.

\textbf{2. Đánh số chữ cái.} Đánh số A $=0$, B $=1$, \ldots, Z $=25$. Mã hoá chữ số $x$ với bước $k$:
lấy $x + k$; nếu kết quả $\ge 26$ thì trừ bớt $26$ (tức là lấy \emph{số dư} khi chia $x+k$ cho $26$).
\begin{center}
\renewcommand{\arraystretch}{1.1}\setlength{\tabcolsep}{2.6pt}\small
\begin{tabular}{|*{13}{c|}}
\hline
\rowcolor{tmBlue!10}A&B&C&D&E&F&G&H&I&J&K&L&M\\ \hline
0&1&2&3&4&5&6&7&8&9&10&11&12\\ \hline
\rowcolor{tmBlue!10}N&O&P&Q&R&S&T&U&V&W&X&Y&Z\\ \hline
13&14&15&16&17&18&19&20&21&22&23&24&25\\ \hline
\end{tabular}
\end{center}
Ví dụ, bước $4$: Y $= 24$, $24 + 4 = 28 \ge 26$, $28 - 26 = 2 =$ C. Vậy Y $\to$ C.

\textbf{3. Dịch nhiều lần.} Dịch $a$ bước rồi dịch tiếp $b$ bước giống như dịch một lần $a + b$ bước. Dịch $26$ bước thì mỗi
chữ quay về chính nó, nên chỉ cần quan tâm đến \emph{số dư} của số bước khi chia cho $26$.

\textbf{4. Ngày trong tuần cũng ``quay vòng''.} Sau Chủ nhật lại đến thứ Hai: các ngày lặp lại theo chu kì $7$. ``$n$ ngày nữa''
là thứ mấy chỉ phụ thuộc vào số dư của $n$ khi chia cho $7$.
\end{knowbox}
\begin{figC}
  \node at (0,0){\lgaWheel{4}};
  \node[align=left,font=\small,anchor=west] at (2.5,0){Bước dịch $4$:\\ A $\to$ E, \ B $\to$ F, \ldots\\ W $\to$ A, \ X $\to$ B,\\ Y $\to$ C, \ Z $\to$ D.};
\end{figC}
\begin{proofbox}
Vì sao dịch $26$ bước thì về chỗ cũ? Vòng tròn có đúng $26$ chữ cái, nên đi $26$ bước là đi trọn một vòng. Tương tự, dịch
$30$ bước $= 26 + 4$ bước, tức là đi trọn một vòng rồi đi thêm $4$ bước: giống hệt dịch $4$ bước. Và $4$ chính là số dư
khi chia $30$ cho $26$.
\end{proofbox}
\begin{warnbox}
Lỗi hay gặp nhất là quên ``quay vòng'' ở cuối bảng chữ cái: W dịch $4$ bước \emph{không} ra một kí tự lạ sau Z, mà
ra A (W $\to$ X $\to$ Y $\to$ Z $\to$ A). Lỗi thứ hai: giải mã lại dịch \emph{xuôi} thay vì dịch \emph{ngược}.
\end{warnbox}

% ---------------------------------------------------------------------
\needspace{9\baselineskip}\shead{Ma phương}
\begin{knowbox}
\textbf{1. Định nghĩa.} Ma phương $3\times 3$ là một bảng $9$ ô, điền $9$ số khác nhau sao cho tổng các số trên mỗi hàng,
mỗi cột và mỗi đường chéo đều bằng nhau. Tổng chung ấy gọi là \textbf{tổng ma phương}, kí hiệu $S$.

\textbf{2. Tổng ma phương.} Ba hàng phủ kín cả bảng, mỗi ô đúng một lần, nên
\[ 3S = \text{tổng của cả } 9 \text{ số}, \qquad S = \frac{\text{tổng của cả } 9 \text{ số}}{3}. \]
Ví dụ: với các số $1, 2, \ldots, 9$: tổng là $45$, nên $S = 15$.

\textbf{3. Số ở ô chính giữa.} Với ma phương $3\times 3$: \ \textbf{số ở giữa $= S : 3$} (bằng trung bình cộng của $9$ số).
Với các số $1, \ldots, 9$, số ở giữa là $5$.

\textbf{4. Mẹo điền nhanh.} Khi đã biết $S$ và có một hàng (cột, đường chéo) đã có hai số, số còn lại $= S -$ (tổng hai số đã biết).
\end{knowbox}
\begin{figC}
  \node at (0,0){\lgaMagic{$2$}{$7$}{$6$}{$9$}{$5$}{$1$}{$4$}{$3$}{$8$}};
  \begin{scope}[xshift=4.2cm]
  \fill[tmOrange!25] (-1.35,-0.45) rectangle (1.35,0.45);
  \fill[tmOrange!25] (-0.45,-1.35) rectangle (0.45,1.35);
  \draw[tmPurple,line width=1.4pt,opacity=0.7] (-1.25,1.25)--(1.25,-1.25) (-1.25,-1.25)--(1.25,1.25);
  \draw[Navy,line width=0.9pt] (-1.35,-1.35) grid[step=0.9] (1.35,1.35);
  \node[font=\bfseries] at (0,0){$c$};
  \node[font=\footnotesize,text=Slate,align=center] at (0,-1.75){$4$ đường đi qua ô giữa};
  \end{scope}
\end{figC}
\begin{proofbox}
Vì sao số ở giữa bằng $S : 3$? Có đúng $4$ đường đi qua ô giữa: hàng giữa, cột giữa và hai đường chéo (hình bên phải).
Cộng tổng của $4$ đường này, ta được $4S$. Trong phép cộng đó, ô giữa $c$ được tính $4$ lần, còn mỗi ô khác được tính
đúng $1$ lần (hãy kiểm tra trên hình: mỗi ô ngoài thuộc đúng một trong bốn đường). Do đó
\[ 4S = (\text{tổng } 9 \text{ số}) + 3c = 3S + 3c. \]
Bớt $3S$ ở cả hai vế: $S = 3c$, tức là $c = S : 3$.
\end{proofbox}
\begin{tipbox}
Ma phương có thể dùng cả số âm. Với các số $-6, -5, \ldots, 2$: tổng là $-18$, nên $S = -6$ và số ở giữa là $-2$. Công thức
vẫn y nguyên!
\end{tipbox}

% ---------------------------------------------------------------------
\needspace{9\baselineskip}\shead{Chiếc cân chính là một phương trình}
\begin{knowbox}
\textbf{1. Cân thăng bằng $=$ đẳng thức.} Nếu chiếc cân thăng bằng thì tổng khối lượng hai đĩa bằng nhau. Gọi khối lượng mỗi
loại vật là một chữ, ta viết được một đẳng thức (phương trình).

\textbf{2. Bỏ bớt vật giống nhau.} Lấy ra khỏi \emph{cả hai} đĩa những vật giống nhau thì cân vẫn thăng bằng
(trừ hai vế cho cùng một số).

\textbf{3. Thay thế.} Nếu một nhóm vật $A$ cân bằng với nhóm $B$, thì ở bất kì chiếc cân nào ta cũng có thể thay $A$ bằng $B$
mà cân vẫn thăng bằng.

\textbf{4. Đổi ra ``đơn vị nhỏ''.} Nếu $4$ táo nặng bằng $3$ cam, hãy coi mỗi quả táo nặng $3$ phần và mỗi quả cam nặng $4$ phần
(vì $4\times 3 = 3\times 4 = 12$ phần). Khi mọi vật đều đổi ra ``phần'', so sánh hai nhóm vật chỉ còn là so sánh hai số.
\end{knowbox}
\begin{figC}
  \node at (0,0){\lgaScale{\lgaApple\,\lgaApple\,\lgaApple\,\lgaApple}{\lgaOrg\,\lgaOrg\,\lgaOrg}};
  \node[font=\small,anchor=west,align=left] at (2.7,0){$4\cdot T = 3\cdot C$\\ $\Rightarrow$ táo $= 3$ phần, cam $= 4$ phần};
\end{figC}
\begin{proofbox}
Vì sao cách ``đổi ra phần'' đúng? Nếu táo nặng $3$ phần, cam nặng $4$ phần thì $4$ táo nặng $12$ phần và $3$ cam cũng nặng
$12$ phần: cân thăng bằng. Hơn nữa, từ $4T = 3C$ suy ra $T = \dfrac{3}{4}C$, nghĩa là tỉ số khối lượng táo : cam \emph{luôn}
bằng $3 : 4$ -- dù quả cam thật nặng bao nhiêu gam đi nữa. Vì chỉ có tỉ số là quan trọng, ta được phép chọn ``$1$ phần''
làm đơn vị.
\end{proofbox}
\begin{warnbox}
Không được bỏ một vật ở \emph{một} đĩa mà không bỏ vật giống hệt ở đĩa kia. Và chỉ được bỏ những vật \emph{giống nhau}
(cùng khối lượng).
\end{warnbox}

% ---------------------------------------------------------------------
\needspace{9\baselineskip}\shead{Đếm khéo: tổng trên hình, xúc xắc, trò chơi công bằng}
\begin{knowbox}
\textbf{1. Mỗi số được đếm mấy lần?} Khi các số được điền vào các vòng tròn trên cạnh một hình và mọi cạnh có cùng tổng $S$,
hãy cộng tổng của \emph{tất cả các cạnh}. Số nằm ở \emph{đỉnh} thuộc hai cạnh nên được đếm hai lần; số nằm \emph{giữa cạnh}
được đếm một lần. Với tam giác:
\[ 3S = (\text{tổng mọi số}) + (\text{tổng ba số ở đỉnh}). \]

\textbf{2. Giới hạn trên -- dưới.} Tổng ba số ở đỉnh không thể nhỏ hơn tổng ba số nhỏ nhất, cũng không thể lớn hơn tổng ba số lớn
nhất. Từ đó ta khoanh vùng được $S$, rồi kiểm tra từng giá trị.

\textbf{3. Xúc xắc.} Trên con xúc xắc chuẩn, \emph{hai mặt đối diện có tổng bằng $7$}: $1$ -- $6$, $2$ -- $5$, $3$ -- $4$.
Tổng sáu mặt là $1+2+3+4+5+6 = 21$. Vì thế: tổng các mặt \emph{không nhìn thấy} $= 21 -$ tổng các mặt nhìn thấy.

\textbf{4. Trò chơi công bằng.} Gieo một con xúc xắc có $6$ kết quả đồng khả năng. Trò chơi ``thắng nếu \ldots'' là
\emph{công bằng} khi khả năng thắng bằng khả năng thua, tức là có đúng $3$ kết quả thắng và $3$ kết quả thua.
\end{knowbox}
\begin{figC}
  \node at (0,0){\lgaTriSix{$a$}{$b$}{$c$}{$x$}{$y$}{$z$}};
  \node[font=\small,anchor=west,align=left] at (1.8,0.1){cạnh trái: $a + x + b$\\ đáy: \hspace{3.2mm}$b + y + c$\\ cạnh phải: $c + z + a$\\[2pt] \textcolor{tmBlue!80!black}{$a, b, c$ xuất hiện $2$ lần.}};
  \node at (7.4,0){\lgaCube{3}{6}{2}};
  \node[font=\footnotesize,text=Slate,align=center] at (7.6,-1.0){thấy $3 + 6 + 2 = 11$\\ không thấy $21 - 11 = 10$};
\end{figC}
\begin{proofbox}
Vì sao $3S = (\text{tổng mọi số}) + (a + b + c)$? Cộng ba cạnh:
$(a + x + b) + (b + y + c) + (c + z + a) = (a + b + c + x + y + z) + (a + b + c)$.
Mỗi số giữa cạnh $x, y, z$ có mặt một lần; mỗi số ở đỉnh $a, b, c$ có mặt hai lần, nên phải cộng thêm $a + b + c$ một lần nữa.
\end{proofbox}
\begin{warnbox}
Với con xúc xắc đặt trên bàn, mặt \emph{đáy} cũng là mặt không nhìn thấy. Đừng quên nó! Cách an toàn nhất là dùng
``tổng $6$ mặt $= 21$'' rồi trừ đi các mặt nhìn thấy.
\end{warnbox}

% ---------------------------------------------------------------------
\needspace{9\baselineskip}\shead{Hình vuông Latin và sudoku}
\begin{knowbox}
\textbf{1. Hình vuông Latin $4\times 4$.} Điền các chữ số $1, 2, 3, 4$ vào lưới $4\times 4$ sao cho mỗi hàng và mỗi cột có đủ
bốn chữ số (không lặp). \textbf{Sudoku $4\times 4$} thêm điều kiện: mỗi ô vuông $2\times 2$ (vẽ viền đậm) cũng có đủ bốn chữ số.

\textbf{2. Ba ``vùng''.} Mỗi ô thuộc đúng ba vùng: một hàng, một cột, một ô $2\times 2$. Chữ số điền vào ô phải \emph{khác}
mọi chữ số đã có trong cả ba vùng đó.

\textbf{3. Hai cách suy luận.}
\begin{itemize}[leftmargin=1.4em,itemsep=1pt,topsep=2pt]
  \item \emph{Loại trừ ở một ô:} gạch đi mọi chữ số đã có trong hàng, cột, ô $2\times 2$ của ô đó; nếu còn đúng một chữ số thì
  điền nó.
  \item \emph{Tìm chỗ cho một chữ số:} trong một vùng, chữ số $d$ chỉ còn đúng một ô có thể đặt thì $d$ phải nằm ở ô đó.
\end{itemize}

\textbf{4. Ghi ``ứng viên''.} Nếu một ô còn hai khả năng, hãy ghi nhỏ cả hai (ví dụ $\{1;4\}$) rồi đi tiếp; thông tin ở chỗ khác
sẽ quyết định giúp.
\end{knowbox}
\begin{figC}
  \node at (0,0){\lgaFour{1/1/$1$,3/1/$3$,2/2/$2$,4/4/$4$}};
  \node[font=\small,anchor=west,align=left] at (1.8,0){Ô hàng $1$, cột $2$: hàng $1$ đã có $1, 3$;\\ ô $2\times 2$ đã có $1, 2$.\\ Còn lại chỉ có thể là $4$.};
\end{figC}
\begin{warnbox}
Đừng ``đoán'' rồi điền luôn. Mỗi chữ số điền vào phải có lí do (bị loại hết các khả năng khác). Nếu buộc phải thử, hãy ghi rõ
``giả sử'' và sẵn sàng quay lại khi gặp mâu thuẫn.
\end{warnbox}
\begin{parentbox}
Câu đố logic rất hợp để chơi cùng con: hãy đọc to đề bài, để con kể lại bằng lời của mình, rồi cùng vẽ sơ đồ hoặc kẻ bảng.
Khi con đưa ra đáp án, đừng nói ngay đúng hay sai, mà hỏi: ``Con kiểm tra lại với \emph{từng} điều kiện của đề chưa?''.
Thói quen tự kiểm tra quý hơn bất kì mẹo giải nào.
\end{parentbox}

% =====================================================================
% VÍ DỤ MẪU
% =====================================================================
\begin{exbox}{Ví dụ 1 \ (Người nói thật -- người nói dối)}
Trên một hòn đảo, mỗi người hoặc luôn nói thật, hoặc luôn nói dối. Em gặp hai người Phong và Quyên.
Phong nói: ``Quyên là người nói dối.'' Quyên nói: ``Phong và tôi là cùng một loại người.''
Ai nói thật, ai nói dối?
\guide
\step{1}{Xét trường hợp Phong nói thật}
Khi đó câu của Phong đúng: Quyên nói dối. Vậy câu của Quyên phải sai. Quyên nói ``hai chúng tôi cùng loại''; thực tế Phong thật,
Quyên dối -- \emph{khác} loại, nên câu của Quyên đúng là sai. \textbf{Không có mâu thuẫn.}
\step{2}{Xét trường hợp Phong nói dối}
Khi đó câu của Phong sai: Quyên \emph{không} nói dối, tức Quyên nói thật. Vậy câu của Quyên phải đúng: ``Phong và Quyên cùng loại''.
Nhưng Phong dối, Quyên thật -- khác loại. Câu của Quyên lại sai. \textbf{Mâu thuẫn!} Trường hợp này bị loại.
\step{3}{Tóm tắt bằng bảng}
\begin{center}
\begin{tabular}{|c|c|c|c|c|}
\hline
\rowcolor{tmBlue!10} Phong & Quyên & Câu của Phong & Câu của Quyên & Kết luận \\ \hline
thật & dối & đúng (hợp lí) & sai (hợp lí) & \textcolor{tmGreen!70!black}{\textbf{được}} \\ \hline
dối & thật & sai (hợp lí) & sai (vô lí!) & \textcolor{AnsRed}{loại} \\ \hline
\end{tabular}
\end{center}
(Hai trường hợp ``cùng thật'' hoặc ``cùng dối'' bị loại ngay từ câu của Phong: nếu Phong thật thì Quyên phải dối; nếu Phong dối thì Quyên phải thật.)
\ans{Phong nói thật, Quyên nói dối.}
\end{exbox}

\begin{exbox}{Ví dụ 2 \ (Mỗi người đoán đúng một nửa)}
Ba đội Đỏ, Xanh, Vàng thi kéo co, không có hai đội nào cùng hạng. Trước trận, ba cổ động viên dự đoán:
\begin{itemize}[leftmargin=1.4em,itemsep=0pt,topsep=2pt]
\item Hằng: ``Đỏ nhất, Xanh nhì.'' \quad Khôi: ``Xanh nhất, Vàng nhì.'' \quad Lâm: ``Xanh ba, Vàng nhất.''
\end{itemize}
Kết quả là mỗi người đoán đúng đúng một ý. Hãy tìm thứ hạng của ba đội.
\guide
\step{1}{Chọn một ý để xét trường hợp}
Mỗi câu có đúng một ý đúng. Ta xét câu của Hằng: hoặc ``Đỏ nhất'' đúng, hoặc ``Xanh nhì'' đúng.
\step{2}{Trường hợp 1: ``Xanh nhì'' đúng (và ``Đỏ nhất'' sai)}
Xanh nhì thì Xanh không nhất, nên ý ``Xanh nhất'' của Khôi sai; vậy ý còn lại ``Vàng nhì'' phải đúng. Nhưng Xanh đã nhì, Vàng
không thể cũng nhì. \textbf{Mâu thuẫn}, loại.
\step{3}{Trường hợp 2: ``Đỏ nhất'' đúng (và ``Xanh nhì'' sai)}
Xanh không nhất (Đỏ nhất), không nhì, nên \textbf{Xanh ba}. Còn lại Vàng nhì.
\step{4}{Thử lại với cả ba người}
\begin{center}
\begin{tabular}{|l|c|c|}
\hline
\rowcolor{tmBlue!10} Người & Ý thứ nhất & Ý thứ hai \\ \hline
Hằng & Đỏ nhất: \textcolor{tmGreen!70!black}{đúng} & Xanh nhì: \textcolor{AnsRed}{sai} \\ \hline
Khôi & Xanh nhất: \textcolor{AnsRed}{sai} & Vàng nhì: \textcolor{tmGreen!70!black}{đúng} \\ \hline
Lâm & Xanh ba: \textcolor{tmGreen!70!black}{đúng} & Vàng nhất: \textcolor{AnsRed}{sai} \\ \hline
\end{tabular}
\end{center}
Mỗi người đúng đúng một ý.
\ans{Đỏ nhất, Vàng nhì, Xanh ba.}
\end{exbox}

\begin{exbox}{Ví dụ 3 \ (Sắp thứ tự -- chắc chắn hay có thể?)}
Có bốn cây bút chì màu đỏ, xanh, vàng, tím, dài khác nhau. Biết: bút đỏ dài hơn bút xanh; bút vàng ngắn hơn bút xanh; bút tím
dài hơn bút vàng nhưng ngắn hơn bút đỏ.\\
a) Bút nào chắc chắn dài nhất? Bút nào chắc chắn ngắn nhất?\\
b) Câu ``Bút tím dài thứ hai'' chắc chắn đúng, có thể đúng hay chắc chắn sai?
\guide
\step{1}{Kí hiệu hoá}
Gọi độ dài là $Đ, X, V, T$. Đề cho: $Đ > X$; \ $X > V$; \ $T > V$; \ $Đ > T$.
\step{2}{Vẽ sơ đồ mũi tên (mũi tên chỉ từ bút ngắn hơn đến bút dài hơn)}
\begin{figC}
  \node[draw=tmOrange,rounded corners=3pt,line width=0.8pt] (V) at (0,0){Vàng};
  \node[draw=tmOrange,rounded corners=3pt,line width=0.8pt] (X) at (2.2,0.7){Xanh};
  \node[draw=tmOrange,rounded corners=3pt,line width=0.8pt] (T) at (2.2,-0.7){Tím};
  \node[draw=tmOrange,rounded corners=3pt,line width=0.8pt] (D) at (4.4,0){Đỏ};
  \draw[->,Navy,line width=0.9pt] (V)--(X); \draw[->,Navy,line width=0.9pt] (V)--(T);
  \draw[->,Navy,line width=0.9pt] (X)--(D); \draw[->,Navy,line width=0.9pt] (T)--(D);
  \draw[<->,dashed,AnsRed!70] (X)--(T) node[midway,right,font=\scriptsize]{?};
\end{figC}
\step{3}{Trả lời câu a}
Từ $Đ > X > V$ và $Đ > T > V$: bút đỏ dài hơn cả ba bút kia, bút vàng ngắn hơn cả ba bút kia. Vậy bút \textbf{đỏ chắc chắn dài nhất},
bút \textbf{vàng chắc chắn ngắn nhất}.
\step{4}{Trả lời câu b: tìm hai trường hợp}
Không có điều kiện nào so sánh Xanh với Tím. Thử hai bộ độ dài (cm) đều thoả mọi điều kiện:
\begin{itemize}[leftmargin=1.4em,itemsep=0pt,topsep=1pt]
\item Đỏ $15$, Tím $13$, Xanh $12$, Vàng $8$: tím dài thứ hai -- câu \emph{đúng};
\item Đỏ $15$, Xanh $12$, Tím $10$, Vàng $8$: tím dài thứ ba -- câu \emph{sai}.
\end{itemize}
Vậy câu ``Bút tím dài thứ hai'' chỉ \textbf{có thể đúng}, không chắc chắn đúng.
\ans{a) Đỏ dài nhất, vàng ngắn nhất (chắc chắn). \ b) Có thể đúng (không chắc chắn).}
\end{exbox}

\begin{exbox}{Ví dụ 4 \ (Mật mã dịch vòng -- tìm bước dịch)}
Bạn Tùng mã hoá một từ bằng mật mã Caesar và nhận được \textbf{JOHV}. Tùng cho biết: với cách mã hoá này, chữ A biến thành chữ H.\\
a) Tìm từ ban đầu. \ b) Sau khi dịch như trên, phải dịch tiếp ít nhất bao nhiêu bước (theo cùng chiều) để mọi chữ quay về chữ ban đầu?
\guide
\step{1}{Tìm bước dịch}
A $= 0$, H $= 7$. A biến thành H nên bước dịch là $k = 7$.
\step{2}{Giải mã: dịch ngược $7$ bước (trừ $7$; nếu ra số âm thì cộng thêm $26$)}
\begin{center}
\begin{tabular}{|c|c|c|c|c|}
\hline
\rowcolor{tmBlue!10} Chữ mã & J & O & H & V \\ \hline
Số & $9$ & $14$ & $7$ & $21$ \\ \hline
Trừ $7$ & $2$ & $7$ & $0$ & $14$ \\ \hline
Chữ gốc & \textbf{C} & \textbf{H} & \textbf{A} & \textbf{O} \\ \hline
\end{tabular}
\end{center}
Từ ban đầu là \textbf{CHAO} (``chào'').
\step{3}{Câu b}
Dịch tổng cộng $26$ bước thì quay đủ một vòng. Đã dịch $7$ bước, cần thêm $26 - 7 = 19$ bước.
\verify{Mã hoá lại CHAO với bước $7$: C$(2)\to 9 =$ J, H$(7)\to 14 =$ O, A$(0)\to 7 =$ H, O$(14)\to 21 =$ V. Được JOHV.}
\ans{a) CHAO. \ b) $19$ bước.}
\end{exbox}

\begin{exbox}{Ví dụ 5 \ (Hoàn thành ma phương)}
Điền các số chẵn $2, 4, 6, \ldots, 18$ (mỗi số một lần) vào bảng sao cho mọi hàng, cột, đường chéo có cùng tổng.
\begin{center}\lgaMagic{$4$}{}{$12$}{}{}{}{}{$6$}{}\end{center}
\guide
\step{1}{Tìm tổng ma phương $S$}
Tổng chín số: $2 + 4 + \cdots + 18 = 2\times(1 + 2 + \cdots + 9) = 2\times 45 = 90$. Vậy $S = 90 : 3 = 30$.
\step{2}{Số ở giữa}
\begin{recallbox}Trong ma phương $3\times 3$, số ở giữa bằng $S : 3$.\end{recallbox}
Số ở giữa $= 30 : 3 = 10$.
\step{3}{Điền dần những đường đã có hai số}
\begin{itemize}[leftmargin=1.4em,itemsep=0pt,topsep=1pt]
\item Hàng trên: ô giữa $= 30 - 4 - 12 = 14$.
\item Đường chéo $4$ -- $10$ -- ?: góc dưới phải $= 30 - 4 - 10 = 16$.
\item Đường chéo $12$ -- $10$ -- ?: góc dưới trái $= 30 - 12 - 10 = 8$.
\item Cột trái: $30 - 4 - 8 = 18$; \ cột phải: $30 - 12 - 16 = 2$.
\end{itemize}
\begin{center}\lgaMagic{$4$}{\textcolor{tmGreen!60!black}{$14$}}{$12$}{\textcolor{tmGreen!60!black}{$18$}}{\textcolor{AnsRed}{$10$}}{\textcolor{tmGreen!60!black}{$2$}}{\textcolor{tmGreen!60!black}{$8$}}{$6$}{\textcolor{tmGreen!60!black}{$16$}}\end{center}
\verify{Hàng dưới: $8 + 6 + 16 = 30$; hàng giữa: $18 + 10 + 2 = 30$; cột giữa: $14 + 10 + 6 = 30$. Mọi số chẵn từ $2$ đến $18$ dùng đúng một lần.}
\ans{Hàng trên $4, 14, 12$; hàng giữa $18, 10, 2$; hàng dưới $8, 6, 16$.}
\end{exbox}

\begin{exbox}{Ví dụ 6 \ (Bài toán chiếc cân)}
Hai chiếc cân dưới đây đều thăng bằng; các hình giống nhau nặng bằng nhau. Một hình vuông nặng bằng bao nhiêu hình tròn?
\begin{center}
\lgaScale{\lgaSq\,\lgaSq}{\lgaTri\,\lgaTri\,\lgaTri}\hspace{1cm}\lgaScale{\lgaTri}{\lgaCir\,\lgaCir}
\end{center}
\guide
\step{1}{Viết thành đẳng thức}
Gọi khối lượng hình vuông, tam giác, tròn là $V, T, O$. Cân 1: $2V = 3T$. \ Cân 2: $T = 2O$.
\step{2}{Thay thế}
Ở cân 1, thay mỗi tam giác bằng hai hình tròn: $2V = 3\times 2O = 6O$.
\step{3}{Chia đôi hai đĩa}
$2V = 6O$ nên $V = 3O$.
\verify{Chọn $O = 1$, $T = 2$, $V = 3$: cân 1: $2\times 3 = 6 = 3\times 2$; cân 2: $2 = 2\times 1$. Đúng.}
\ans{Một hình vuông nặng bằng $3$ hình tròn.}
\end{exbox}

\begin{exbox}{Ví dụ 7 \ (Tổng trên các cạnh tam giác)}
Điền các số $1, 2, 3, 4, 5, 6$ vào sáu vòng tròn (ba ở đỉnh, ba ở giữa cạnh) sao cho tổng ba số trên mỗi cạnh đều bằng $S$.
Tìm tất cả các giá trị có thể của $S$.
\guide
\step{1}{Đếm mỗi số được tính mấy lần}
\begin{recallbox}Cộng tổng ba cạnh: số ở đỉnh được tính hai lần, số ở giữa cạnh được tính một lần.\end{recallbox}
$3S = (1 + 2 + \cdots + 6) + (\text{tổng ba đỉnh}) = 21 + (\text{tổng ba đỉnh})$.
\step{2}{Khoanh vùng}
Tổng ba đỉnh nhỏ nhất là $1 + 2 + 3 = 6$, lớn nhất là $4 + 5 + 6 = 15$. Vậy $27 \le 3S \le 36$, tức $9 \le S \le 12$.
\step{3}{Kiểm tra từng giá trị bằng một cách điền}
\begin{figC}
  \node at (0,0){\lgaTriSix{$1$}{$2$}{$3$}{$6$}{$4$}{$5$}}; \node[font=\small] at (0,-1.6){$S = 9$};
  \node at (3.3,0){\lgaTriSix{$1$}{$3$}{$5$}{$6$}{$2$}{$4$}}; \node[font=\small] at (3.3,-1.6){$S = 10$};
  \node at (6.6,0){\lgaTriSix{$2$}{$4$}{$6$}{$5$}{$1$}{$3$}}; \node[font=\small] at (6.6,-1.6){$S = 11$};
  \node at (9.9,0){\lgaTriSix{$4$}{$5$}{$6$}{$3$}{$1$}{$2$}}; \node[font=\small] at (9.9,-1.6){$S = 12$};
\end{figC}
Chẳng hạn với $S = 10$: $1 + 6 + 3 = 10$, \ $3 + 2 + 5 = 10$, \ $5 + 4 + 1 = 10$.
\ans{$S$ có thể là $9$, $10$, $11$ hoặc $12$.}
\end{exbox}

\begin{exbox}{Ví dụ 8 \ (Sudoku $4\times 4$)}
Hoàn thành lưới sudoku: mỗi hàng, mỗi cột và mỗi ô $2\times 2$ viền đậm đều có đủ các chữ số $1, 2, 3, 4$.
\begin{center}\lgaFour{1/1/$2$,2/1/$1$,3/2/$1$,2/3/$4$,4/4/$3$}\end{center}
\guide
\step{1}{Ô $2\times 2$ trên trái}
Đã có $2, 1$ ở hàng $1$, nên hàng $2$ cột $1$ và cột $2$ là $\{3; 4\}$. Cột $2$ đã có $4$ (ở hàng $3$), nên
hàng $2$ cột $2$ là $3$, hàng $2$ cột $1$ là $4$.
\step{2}{Hàng $2$ và ô trên phải}
Hàng $2$: $4, 3, 1, \_$ $\Rightarrow$ ô cuối là $2$. Hàng $1$ cột $3$, cột $4$ là $\{3; 4\}$; cột $4$ đã có $3$ (hàng $4$) nên
hàng $1$ cột $4$ là $4$, hàng $1$ cột $3$ là $3$.
\step{3}{Cột $2$ và ô dưới trái}
Cột $2$: $1, 3, 4, \_$ $\Rightarrow$ hàng $4$ cột $2$ là $2$. Cột $1$ còn thiếu $\{1; 3\}$; hàng $4$ đã có $3$ nên hàng $4$ cột $1$ là $1$,
hàng $3$ cột $1$ là $3$.
\step{4}{Hai hàng cuối}
Hàng $3$: $3, 4, \_, \_$ còn $\{1; 2\}$; cột $3$ đã có $1$ (hàng $2$) nên hàng $3$ cột $3$ là $2$, cột $4$ là $1$. Hàng $4$: $1, 2, \_, 3$ $\Rightarrow$ ô trống là $4$.
\begin{center}\lgaFour{1/1/$2$,2/1/$1$,3/1/\textcolor{tmGreen!60!black}{$3$},4/1/\textcolor{tmGreen!60!black}{$4$},1/2/\textcolor{tmGreen!60!black}{$4$},2/2/\textcolor{tmGreen!60!black}{$3$},3/2/$1$,4/2/\textcolor{tmGreen!60!black}{$2$},1/3/\textcolor{tmGreen!60!black}{$3$},2/3/$4$,3/3/\textcolor{tmGreen!60!black}{$2$},4/3/\textcolor{tmGreen!60!black}{$1$},1/4/\textcolor{tmGreen!60!black}{$1$},2/4/\textcolor{tmGreen!60!black}{$2$},3/4/\textcolor{tmGreen!60!black}{$4$},4/4/$3$}\end{center}
\verify{Mỗi hàng, mỗi cột, mỗi ô $2\times 2$ đều có đủ $1, 2, 3, 4$.}
\ans{Hàng $1$: $2\,1\,3\,4$; hàng $2$: $4\,3\,1\,2$; hàng $3$: $3\,4\,2\,1$; hàng $4$: $1\,2\,4\,3$.}
\end{exbox}

\clearpage
\dividerbar{Luyện tập cơ bản}

% ---------------------------------------------------------------- Bài 1
\begin{probbox}{1}{\relax}
Mẹ vừa mua một chậu hoa mới. Ba bạn đoán xem đó là hoa gì. An đoán: ``hoa hồng màu đỏ''. Bình đoán: ``hoa cúc màu tím''.
Chi đoán: ``hoa thược dược màu đỏ''. Mẹ cho biết: mỗi bạn đoán đúng \emph{đúng một} trong hai điều -- hoặc màu hoa, hoặc loại hoa.
Chậu hoa mẹ mua là
\choicesII{thược dược tím}{hồng tím}{thược dược đỏ}{hồng vàng}{cúc đỏ}
\end{probbox}
\probsol{
\textbf{Phân tích đề.} Mỗi lời đoán gồm hai ý: \emph{màu} và \emph{loại hoa}. Mỗi bạn đúng đúng một ý: ý này đúng thì ý kia sai.
Cần tìm màu và loại hoa thật.
\begin{recallbox}Câu ``đúng một nửa'': nếu biết một ý đúng thì ý còn lại sai, và ngược lại.\end{recallbox}
\step{1}{Xét màu hoa: có phải màu đỏ không?}
\emph{Giả sử hoa không đỏ.} Khi đó ý ``màu đỏ'' của An sai, nên ý còn lại ``hoa hồng'' phải đúng. Cũng vậy, ý ``màu đỏ'' của Chi sai,
nên ``hoa thược dược'' phải đúng. Hoa vừa là hồng vừa là thược dược -- \textbf{vô lí}. Vậy \textbf{hoa màu đỏ}.
\step{2}{Suy ra loại hoa}
Hoa màu đỏ nên An đúng ý màu $\Rightarrow$ ý ``hoa hồng'' của An sai: hoa không phải hoa hồng.\\
Chi cũng đúng ý màu $\Rightarrow$ ý ``thược dược'' của Chi sai: hoa không phải thược dược.\\
Bình đoán ``màu tím'' -- sai (hoa đỏ), nên ý ``hoa cúc'' của Bình phải đúng: \textbf{hoa cúc}.
\step{3}{Thử lại bằng bảng}
\begin{center}
\begin{tabular}{|l|c|c|c|}
\hline
\rowcolor{tmBlue!10} Bạn & Màu đoán & Loại đoán & Số ý đúng (hoa thật: cúc đỏ) \\ \hline
An & đỏ \textcolor{tmGreen!70!black}{(đúng)} & hồng \textcolor{AnsRed}{(sai)} & $1$ \\ \hline
Bình & tím \textcolor{AnsRed}{(sai)} & cúc \textcolor{tmGreen!70!black}{(đúng)} & $1$ \\ \hline
Chi & đỏ \textcolor{tmGreen!70!black}{(đúng)} & thược dược \textcolor{AnsRed}{(sai)} & $1$ \\ \hline
\end{tabular}
\end{center}
\begin{warnbox}Thử ngay phương án (B) ``hồng tím'': An đúng loại (hồng), Bình đúng màu (tím), nhưng Chi \emph{không đúng ý nào}
-- trái với đề. Phải kiểm tra \emph{cả ba} bạn, không chỉ một hai bạn.\end{warnbox}
\verify{Với ``cúc đỏ'' mỗi bạn đúng đúng một ý. Các phương án khác đều làm ít nhất một bạn đúng $0$ hoặc $2$ ý: (A) thược dược tím -- An đúng $0$ ý; (C) thược dược đỏ -- Chi đúng $2$ ý; (D) hồng vàng -- Bình đúng $0$ ý.}
Vậy chọn \textbf{(E)}.
\ans{(E) cúc đỏ}
}

% ---------------------------------------------------------------- Bài 2
\begin{probbox}{2}{\relax}
$26$ chữ cái tiếng Anh được viết theo thứ tự A, B, C, \ldots, Z quanh một vòng tròn, theo chiều kim đồng hồ. Để \emph{mã hoá} một
tin nhắn, ta thay mỗi chữ cái bằng chữ cái đứng sau nó $4$ vị trí theo chiều kim đồng hồ (đây gọi là \emph{mật mã Caesar}).
Ví dụ, từ BOX được mã hoá thành FSB. Từ YEU được mã hoá thành
\choices{CIZ}{BHX}{CIY}{UAQ}{DJZ}
\end{probbox}
\probsol{
\textbf{Phân tích đề.} Bước dịch là $4$, chiều xuôi (chiều kim đồng hồ). Cần mã hoá từng chữ Y, E, U.
\begin{recallbox}Đánh số A $= 0$, \ldots, Z $= 25$. Mã hoá: cộng $4$; nếu được số $\ge 26$ thì trừ $26$ (đi quá Z thì quay về A).\end{recallbox}
\step{1}{Kiểm tra cách hiểu bằng ví dụ BOX $\to$ FSB}
B$(1) + 4 = 5 =$ F; \ O$(14) + 4 = 18 =$ S; \ X$(23) + 4 = 27$, $27 - 26 = 1 =$ B. Đúng là FSB, và ta thấy chữ X ``quay vòng'' thành B.
\step{2}{Mã hoá từng chữ của YEU}
\begin{center}
\begin{tabular}{|c|c|c|c|}
\hline
\rowcolor{tmBlue!10} Chữ gốc & Y & E & U \\ \hline
Số & $24$ & $4$ & $20$ \\ \hline
Cộng $4$ & $28 \to 28 - 26 = 2$ & $8$ & $24$ \\ \hline
Chữ mã & \textbf{C} & \textbf{I} & \textbf{Y} \\ \hline
\end{tabular}
\end{center}
\begin{figC}
  \node at (0,0){\lgaWheel{4}};
  \node[align=left,font=\small,anchor=west] at (2.5,0){Y $\to$ Z $\to$ A $\to$ B $\to$ \textbf{C}\\ E $\to$ F $\to$ G $\to$ H $\to$ \textbf{I}\\ U $\to$ V $\to$ W $\to$ X $\to$ \textbf{Y}};
\end{figC}
\begin{warnbox}(B) BHX là dịch $3$ bước (đếm thiếu một bước); (D) UAQ là dịch \emph{ngược} $4$ bước (giải mã thay vì mã hoá);
(E) DJZ là dịch $5$ bước. Hãy đếm cẩn thận: ``sau $4$ vị trí'' nghĩa là đi $4$ bước, không tính chữ xuất phát.\end{warnbox}
\verify{Giải mã CIY bằng cách lùi $4$ bước: C $\to$ B $\to$ A $\to$ Z $\to$ Y; \ I $\to$ E; \ Y $\to$ U. Được lại YEU.}
Vậy chọn \textbf{(C)}.
\ans{(C) CIY}
}

% ---------------------------------------------------------------- Bài 3
\begin{probbox}{3}{\relax}
Năm bạn Minh, Hùng, Phúc, Tâm, Quân thi chạy, không có hai bạn nào về đích cùng lúc. Minh chạy nhanh hơn Hùng và Phúc. Hùng chạy
chậm hơn Phúc. Tâm chạy nhanh hơn Minh nhưng chậm hơn Quân. Bạn nào về đích thứ tư?
\choices{Phúc}{Hùng}{Quân}{Minh}{Tâm}
\end{probbox}
\probsol{
\textbf{Phân tích đề.} Có năm bạn, các câu so sánh tốc độ. Cần sắp xếp thứ tự về đích và tìm người thứ tư. Ai nhanh hơn thì về trước.
\step{1}{Kí hiệu hoá (dấu $>$ nghĩa là ``nhanh hơn'')}
\begin{itemize}[leftmargin=1.4em,itemsep=0pt,topsep=1pt]
\item ``Minh nhanh hơn Hùng và Phúc'': \ $M > H$ \ và \ $M > P$.
\item ``Hùng chậm hơn Phúc'': \ $P > H$.
\item ``Tâm nhanh hơn Minh nhưng chậm hơn Quân'': \ $Q > T > M$.
\end{itemize}
\step{2}{Nối thành một dãy (tính bắc cầu)}
Từ $Q > T > M$, \ $M > P$ và $P > H$ ta được: \quad $Q > T > M > P > H$.
\begin{figC}
  \foreach \x/\n/\k in {0/Quân/1,2/Tâm/2,4/Minh/3,6/Phúc/4,8/Hùng/5}{
    \node[draw=tmBlue,fill=tmBlue!8,rounded corners=3pt,minimum width=1.3cm,line width=0.8pt] (n\k) at (\x,0){\n};
    \node[font=\scriptsize,text=Slate] at (\x,-0.5){thứ \k};}
  \foreach \a/\b in {1/2,2/3,3/4,4/5}{\draw[->,Navy,line width=0.8pt] (n\a)--(n\b);}
  \node[font=\scriptsize,text=Slate] at (4,0.55){(mũi tên: ``nhanh hơn'')};
  \draw[AnsRed,line width=1.2pt,rounded corners=4pt] (5.25,-0.8) rectangle (6.75,0.35);
\end{figC}
\step{3}{Đọc kết quả}
Thứ tự về đích: Quân (nhất), Tâm (nhì), Minh (ba), \textbf{Phúc (tư)}, Hùng (năm).
\begin{warnbox}Câu ``Minh nhanh hơn Hùng và Phúc'' chưa cho biết Hùng và Phúc ai nhanh hơn -- phải dùng câu ``Hùng chậm hơn Phúc''.
Đừng bỏ qua câu nào!\end{warnbox}
\verify{Minh (ba) nhanh hơn Hùng (năm) và Phúc (tư) \checkmark; Hùng chậm hơn Phúc \checkmark; Tâm (nhì) nhanh hơn Minh, chậm hơn Quân \checkmark.}
Vậy chọn \textbf{(A)}.
\ans{(A) Phúc}
}

% ---------------------------------------------------------------- Bài 4
\begin{probbox}{4}{\relax}
\begin{minipage}[t]{0.62\linewidth}
Trong một ma phương, tổng các số trên mỗi hàng, mỗi cột và mỗi đường chéo đều bằng nhau. Ma phương bên cạnh dùng mỗi số
nguyên từ $-6$ đến $2$ đúng một lần. Giá trị của $Y$ là
\choicesII{$-1$}{$0$}{$-6$}{$2$}{$-2$}
\end{minipage}\hfill
\begin{minipage}[t]{0.33\linewidth}\vspace{-2pt}\centering
\lgaMagic{$1$}{}{$Y$}{$-4$}{}{}{$-3$}{}{$-5$}
\end{minipage}
\end{probbox}
\probsol{
\textbf{Phân tích đề.} Chín số là $-6, -5, -4, -3, -2, -1, 0, 1, 2$. Đã biết bốn ô; cần tìm ô $Y$ ở góc trên bên phải.
\step{1}{Tìm tổng ma phương $S$}
Tổng chín số: $(-6) + (-5) + (-4) + (-3) + (-2) + (-1) + 0 + 1 + 2$. Ghép $(-2) + 2 = 0$, $(-1) + 1 = 0$, còn lại
$(-6) + (-5) + (-4) + (-3) = -18$. Vậy tổng là $-18$, và
\[ S = (-18) : 3 = -6. \]
\verify{Cột trái đã đủ ba số: $1 + (-4) + (-3) = -6$. Khớp với $S = -6$.}
\step{2}{Tìm số ở ô giữa}
\begin{recallbox}Trong ma phương $3\times 3$, số ở giữa bằng $S : 3$.\end{recallbox}
Số ở giữa $= (-6) : 3 = -2$. (Cách khác: đường chéo chính $1 + c + (-5) = -6$ nên $c = -6 - 1 + 5 = -2$.)
\step{3}{Tìm $Y$ bằng đường chéo phụ}
Đường chéo đi qua $Y$, ô giữa và ô $-3$: \ $Y + (-2) + (-3) = -6$, nên $Y = -6 + 2 + 3 = -1$.
\step{4}{Hoàn thành bảng để thử lại}
Hàng trên: ô giữa $= -6 - 1 - (-1) = -6$. Hàng giữa: ô phải $= -6 - (-4) - (-2) = 0$. Hàng dưới: ô giữa $= -6 - (-3) - (-5) = 2$.
\begin{center}\lgaMagic{$1$}{\textcolor{tmGreen!60!black}{$-6$}}{\textcolor{AnsRed}{$-1$}}{$-4$}{\textcolor{tmGreen!60!black}{$-2$}}{\textcolor{tmGreen!60!black}{$0$}}{$-3$}{\textcolor{tmGreen!60!black}{$2$}}{$-5$}\end{center}
\verify{Cột giữa: $-6 - 2 + 2 = -6$; cột phải: $-1 + 0 - 5 = -6$; hàng giữa: $-4 - 2 + 0 = -6$. Mỗi số từ $-6$ đến $2$ dùng đúng một lần.}
\begin{warnbox}Cẩn thận với dấu trừ: $-6 - (-4) = -6 + 4 = -2$. Trừ một số âm là cộng với số đối của nó.\end{warnbox}
Vậy chọn \textbf{(A)}.
\ans{(A) $-1$}
}

% ---------------------------------------------------------------- Bài 5
\begin{probbox}{5}{\relax}
Một trò chơi được gọi là \emph{công bằng} nếu khả năng thắng bằng khả năng thua. Gieo một con xúc xắc $6$ mặt cân đối. Trong
bốn trò chơi sau, có bao nhiêu trò chơi công bằng?
\begin{itemize}[leftmargin=1.6em,itemsep=0pt,topsep=2pt]
\item Em thắng nếu gieo được mặt $2$ chấm.
\item Em thắng nếu gieo được số chẵn.
\item Em thắng nếu gieo được số nhỏ hơn $4$.
\item Em thắng nếu gieo được số chia hết cho $3$.
\end{itemize}
\choices{$0$}{$1$}{$2$}{$3$}{$4$}
\end{probbox}
\probsol{
\textbf{Phân tích đề.} Có $6$ kết quả đồng khả năng: $1, 2, 3, 4, 5, 6$. Trò chơi công bằng khi số kết quả thắng bằng số kết quả thua,
tức là có $3$ kết quả thắng, $3$ kết quả thua.
\begin{recallbox}Gieo xúc xắc cân đối: sáu mặt có khả năng xuất hiện như nhau. Khả năng thắng $=$ (số kết quả thắng) $:$ $6$.\end{recallbox}
\step{1}{Liệt kê kết quả thắng của từng trò chơi}
\begin{center}
\begin{tabular}{|l|c|c|c|c|}
\hline
\rowcolor{tmBlue!10} Trò chơi & Kết quả thắng & Số thắng & Số thua & Công bằng? \\ \hline
Mặt $2$ chấm & $2$ & $1$ & $5$ & không \\ \hline
Số chẵn & $2, 4, 6$ & $3$ & $3$ & \textbf{có} \\ \hline
Nhỏ hơn $4$ & $1, 2, 3$ & $3$ & $3$ & \textbf{có} \\ \hline
Chia hết cho $3$ & $3, 6$ & $2$ & $4$ & không \\ \hline
\end{tabular}
\end{center}
\step{2}{Đếm}
Có $2$ trò chơi công bằng: ``số chẵn'' và ``nhỏ hơn $4$''.
\begin{warnbox}``Nhỏ hơn $4$'' \emph{không} gồm số $4$: chỉ có $1, 2, 3$. Nếu tính nhầm cả $4$ thì có $4$ kết quả thắng -- không công bằng.
Còn ``chia hết cho $3$'' gồm $3$ và $6$ (đừng quên $6$).\end{warnbox}
\verify{Trò ``số chẵn'': khả năng thắng $\dfrac{3}{6} = \dfrac12$, khả năng thua cũng $\dfrac12$. Trò ``nhỏ hơn $4$'' cũng vậy.}
Vậy chọn \textbf{(C)}.
\ans{(C) $2$}
}

% ---------------------------------------------------------------- Bài 6
\begin{probbox}{6}{\relax}
Năm bạn Giang, Khoa, Lan, Mai, Nam ngồi vào $5$ chiếc ghế quanh một chiếc bàn tròn. Lan và Mai ngồi cạnh nhau. Giang và Khoa
\emph{không} ngồi cạnh nhau. Hai bạn ngồi hai bên cạnh Nam là
\choicesII{Giang và Lan}{Giang và Khoa}{Khoa và Mai}{Lan và Mai}{Mai và Giang}
\end{probbox}
\probsol{
\textbf{Phân tích đề.} Năm người quanh bàn tròn $5$ ghế. Có hai điều kiện: Lan -- Mai cạnh nhau; Giang -- Khoa không cạnh nhau. Cần tìm
hai người ngồi cạnh Nam.
\begin{recallbox}Quanh bàn tròn có thể cố định chỗ của một người (xoay bàn không làm thay đổi ai cạnh ai). Với $5$ ghế, mỗi người
không ngồi cạnh đúng $2$ người -- đó là hai người ``cách một ghế''.\end{recallbox}
\step{1}{Bắt đầu từ điều kiện chặt nhất: Giang và Khoa không cạnh nhau}
Đánh số ghế $1$ đến $5$ theo vòng. Đặt Giang ở ghế $1$. Hai ghế cạnh ghế $1$ là ghế $2$ và ghế $5$, nên Khoa ngồi ghế $3$ hoặc ghế $4$.
Hai trường hợp này đối xứng nhau (lật bàn như soi gương), nên ta xét Khoa ở ghế $3$.
\step{2}{Đặt Lan và Mai}
Còn trống ghế $2$, $4$, $5$. Lan và Mai phải ngồi hai ghế \emph{cạnh nhau}. Trong ba ghế này chỉ có ghế $4$ và ghế $5$ cạnh nhau
(ghế $2$ cạnh ghế $1$ và ghế $3$, đều đã có người). Vậy Lan, Mai ngồi ghế $4$ và $5$.
\step{3}{Nam ngồi ghế còn lại}
Nam ngồi ghế $2$, nằm giữa ghế $1$ (Giang) và ghế $3$ (Khoa).
\begin{figC}
  \foreach \i in {1,...,5}{\coordinate (S\i) at ({90-(\i-1)*72}:1.35);}
  \draw[Navy!40,line width=5pt] (0,0) circle (0.75);
  \foreach \i/\n/\c in {1/Giang/tmBlue,2/Nam/AnsRed,3/Khoa/tmBlue,4/Lan/tmGreen,5/Mai/tmGreen}{
    \node[circle,draw=\c,fill=\c!10,line width=0.9pt,minimum size=0.95cm,font=\scriptsize] at (S\i){\n};
    \node[font=\scriptsize,text=Slate] at ({90-(\i-1)*72}:2.05){ghế \i};}
\end{figC}
(Nếu Khoa ở ghế $4$ thì tương tự: Lan, Mai ở ghế $2$, $3$ và Nam ở ghế $5$ -- vẫn ngồi giữa Giang và Khoa.)
\begin{proofbox}Vì sao không còn cách xếp nào khác? Giang và Khoa không cạnh nhau nên giữa họ có đúng một ghế ở một phía và hai ghế
ở phía kia. Lan -- Mai cần hai ghế liền nhau, nên phải chiếm phía có hai ghế. Ghế lẻ loi giữa Giang và Khoa thuộc về Nam.\end{proofbox}
\verify{Lan -- Mai cạnh nhau \checkmark; Giang (ghế $1$) và Khoa (ghế $3$) cách nhau một ghế \checkmark.}
Vậy chọn \textbf{(B)}.
\ans{(B) Giang và Khoa}
}

% ---------------------------------------------------------------- Bài 7
\begin{probbox}{7}{\relax}
\begin{minipage}[t]{0.62\linewidth}
Trên mỗi con xúc xắc, tổng số chấm trên hai mặt đối diện luôn bằng $7$. Hai con xúc xắc được đặt trên bàn như hình bên. Tổng số
chấm trên tất cả các mặt \emph{không nhìn thấy} của hai con xúc xắc là
\choicesII{$14$}{$20$}{$21$}{$24$}{$30$}
\end{minipage}\hfill
\begin{minipage}[t]{0.34\linewidth}\vspace{2pt}\centering
\lgaCube{3}{6}{2}\hspace{0.3cm}\lgaCube{1}{5}{4}
\end{minipage}
\end{probbox}
\probsol{
\textbf{Phân tích đề.} Mỗi con xúc xắc có $6$ mặt, ta nhìn thấy $3$ mặt (trên, trước, bên phải), còn $3$ mặt không nhìn thấy
(dưới, sau, bên trái). Cần tổng số chấm trên $6$ mặt không nhìn thấy của cả hai con.
\begin{recallbox}Tổng số chấm trên sáu mặt của một con xúc xắc là $1 + 2 + 3 + 4 + 5 + 6 = 21$.\end{recallbox}
\step{1}{Đọc các mặt nhìn thấy}
Con thứ nhất: trên $3$, trước $6$, phải $2$ $\Rightarrow$ tổng $3 + 6 + 2 = 11$.\\
Con thứ hai: trên $1$, trước $5$, phải $4$ $\Rightarrow$ tổng $1 + 5 + 4 = 10$.
\step{2}{Tính các mặt không nhìn thấy}
Con thứ nhất: $21 - 11 = 10$. \quad Con thứ hai: $21 - 10 = 11$. \quad Tổng: $10 + 11 = 21$.
\step{3}{Cách thứ hai: dùng mặt đối diện}
Mỗi mặt không nhìn thấy đối diện với một mặt nhìn thấy, nên nó có $7 - (\text{mặt nhìn thấy})$ chấm.
Con thứ nhất: $(7 - 3) + (7 - 6) + (7 - 2) = 4 + 1 + 5 = 10$. Con thứ hai: $(7 - 1) + (7 - 5) + (7 - 4) = 6 + 2 + 3 = 11$.
\begin{proofbox}Vì sao mỗi mặt không thấy đối diện đúng một mặt thấy? Ba mặt thấy (trên, trước, phải) đôi một kề nhau, không có hai
mặt nào đối diện. Ba mặt đối diện với chúng (dưới, sau, trái) chính là ba mặt bị che.\end{proofbox}
\begin{warnbox}Có bạn chỉ tính hai mặt bị che ``phía sau'' mà quên mặt đáy áp xuống bàn. Mặt đáy cũng không nhìn thấy!\end{warnbox}
\verify{Tổng mọi mặt của hai con là $21 + 21 = 42$; mặt thấy $11 + 10 = 21$; mặt không thấy $42 - 21 = 21$.}
Vậy chọn \textbf{(C)}.
\ans{(C) $21$}
}
